  % !Mode:: "TeX:UTF-8"
\documentclass[12pt]{book}
%%%%=========加载宏包=======
\usepackage{pdfpages}
\usepackage{fontspec}
\usepackage{xeCJK}
\setCJKmainfont{FangSong_GB2312}
\setmainfont{TimesNewRomanPSMT}
%\setCJKmainfont{STKaitiSC-Regular}
\usepackage{graphics}           % 插入图片
\usepackage{epstopdf}
\usepackage{color}              % 支持彩色
\usepackage{xcolor,graphicx}
\usepackage{geometry}
\usepackage{float}
\usepackage{indentfirst}        % 首行缩进宏包
\usepackage{amsmath,amssymb,bm} % AMS math 宏包与数学符号加粗
\usepackage{amsthm}
\usepackage{cases}              % 数学公式
\usepackage{setspace}
\usepackage{hyperref}
\usepackage{titlesec}
\usepackage[all]{xy}            % 交换图
\usepackage{tikz}
\usepackage{multirow}% 合并 Table 的单元格时用到
\usepgflibrary{arrows}
\hypersetup{dvips,ps2pdf,CJKbookmarks,hypertex,
	colorlinks=true,
	pdfpagemode=FullScreen,
	bookmarksopen=true,
	bookmarksnumbered=true}     % 在中文环境下设置超链接文本
\usepackage{graphicx} %use graph format
\usepackage{epstopdf}
\usepackage{caption}
\usepackage{diagbox}
\usepackage{mathrsfs}
\usepackage{tikz}
\usepackage{dsfont}
\usepackage{graphics}
\usepackage{amsmath,amssymb,amsfonts,amsthm,bm,mathrsfs}
\usepackage{epsfig,graphicx,epstopdf,picinpar,subfigure,pstricks}
\usepackage{hyperref,cleveref,cite}
\usepackage{titlesec}
\usepackage{tcolorbox}
\usepackage{fancyhdr,fancyvrb}
\usepackage{xcolor,graphicx,geometry,float,indentfirst,cases,setspace,titlesec}
\usepackage{diagbox}



%%%%%%%%%====设置页面=======================%%%%%
\setlength{\paperwidth}{21cm}
\setlength{\paperheight}{29.7cm}
\setlength{\textwidth}{14cm}  % 正文宽度
\setlength{\textheight}{24.0cm}   % 正文高度
\setlength{\hoffset}{0cm}       % 左边距 = \hoffset + 1 英寸
\setlength{\voffset}{0cm}       % 顶端距离 = \voffset + 1 英寸
\setlength{\parskip}{3pt plus1pt minus1pt}  % 段落之间的竖直距离
\renewcommand{\baselinestretch}{1.22}    % 定义行距
\topmargin=-1cm

%%%%%%%%%%%====设置字体===========%%%
\setCJKfamilyfont{songti}{STSongti-SC-Regular} %宋体
\newcommand{\song}{\CJKfamily{songti}}    % 宋体
\setCJKfamilyfont{kaiti}{STKaitiSC-Regular} %楷体
\newcommand{\kaishu}{\CJKfamily{kaiti}}      % 楷体
\setCJKfamilyfont{fangsong}{FangSong_GB2312} %仿宋
\newcommand{\fs}{\CJKfamily{fongsong}}
\setCJKfamilyfont{hei}{STHeitiSC-Medium} %黑体
\newcommand{\hei}{\CJKfamily{hei}}      % 黑体

%%%%========定义数学环境======
\newtheoremstyle{mythm}{3pt}{3pt}{\kaishu}{\parindent}{\hei}{}{2em}{}
\theoremstyle{mythm}
\newtheorem{theorem}{定理}[section]
\newtheorem{lemma}{引理}[section]
\newtheorem{corollary}{推论}[section]
\newtheorem{definition}{定义}[section]
\newtheorem{proposition}{命题}[section]
\newtheorem{property}{性质}[section]
\newtheorem{example}{例}[section]
\newtheorem{remark}{注记}[section]
\newtheorem{conjecture}{猜想}[section]
\renewcommand{\proofname}{\hei{证明}}
%%%设置罗马字体
\newcommand{\NRoman}[1]{\uppercase\expandafter{\romannumeral #1}}
\newcommand{\nroman}[1]{\lowercase\expandafter{\romannumeral #1}}

%%============================设置章节格式===================================%%
\titleformat{\section}[hang]{\hei \erhao \centering }{\erhao \hei \S\,\thesection}{1em}{}{}
\titlespacing{\section}{0pt}{1.5ex plus .1ex minus .2ex}{\wordsep}
\titleformat{\subsection}[hang]{\normalfont\bf}{\thesubsection}{1em}{}{}
\titlespacing{\subsubsection}{2em}{1.5ex plus .1ex minus .2ex}{\wordsep}


%%%%%%%%%%%%%%%%%%%%%%%%%%%%%%%%%%%%%%%%%%%%%%%%%%%%%%%%%%%%%%%%
%%       中文字体设置                                         %%
%%%%%%%%%%%%%%%%%%%%%%%%%%%%%%%%%%%%%%%%%%%%%%%%%%%%%%%%%%%%%%%%

\newcommand{\chuhao}{\fontsize{42pt}{\baselineskip}\selectfont}     % 初号
\newcommand{\xiaochuhao}{\fontsize{36pt}{\baselineskip}\selectfont} % 小初号
\newcommand{\yihao}{\fontsize{28pt}{\baselineskip}\selectfont}      % 一号
\newcommand{\xiaoyihao}{\fontsize{24pt}{\baselineskip}\selectfont}  % 小一号
\newcommand{\erhao}{\fontsize{21pt}{\baselineskip}\selectfont}      % 二号
\newcommand{\xiaoerhao}{\fontsize{18pt}{\baselineskip}\selectfont}  % 小二号
\newcommand{\sanhao}{\fontsize{16pt}{\baselineskip}\selectfont}     % 三号
\newcommand{\xiaosanhao}{\fontsize{15pt}{\baselineskip}\selectfont} % 小三号
\newcommand{\sihao}{\fontsize{14pt}{\baselineskip}\selectfont}      % 四号
\newcommand{\xiaosihao}{\fontsize{12pt}{\baselineskip}\selectfont}  % 小四号
\newcommand{\wuhao}{\fontsize{10.5pt}{\baselineskip}\selectfont}    % 五号
\newcommand{\xiaowuhao}{\fontsize{9pt}{\baselineskip}\selectfont}   % 小五号



%%%%%%%%%%%%%%%%%%%%%%%%%%%%%%%%%%%%%%%%%%%%%%%%%%%%%%%%%%%%%%%%
%       重新定义页脚线长                                       %
%%%%%%%%%%%%%%%%%%%%%%%%%%%%%%%%%%%%%%%%%%%%%%%%%%%%%%%%%%%%%%%%
\makeatletter
\renewcommand\footnoterule{\kern-3\p@ \hrule width 0.15\columnwidth \kern 2.6\p@}
\makeatother
\renewcommand{\thefootnote}{\fnsymbol{footnote}}




\renewcommand\figurename{\hei 图}    
\renewcommand\tablename{\hei 表}    
\renewcommand{\contentsname}{目~录}


\bibliographystyle{plain}

\allowdisplaybreaks[4]


%%%%%========设置插入图片的路径=======%%%%
%\graphicspath{{Figure/}} 

%%%%%===设置默认路径=========%%%%
%\makeatletter
%\providecommand*\input@path{}
%\newcommand\addinputpath[1]{
%\expandafter\def\expandafter\input@path
%\expandafter{\input@path{#1}}}
%\addinputpath{/Users/dennisjxq/OneDrive - tju.edu.cn/tex/}
%\makeatother


\begin{document}

\title{高数下知识点小结}
\maketitle

\setcounter{chapter}{8}

\chapter{多元函数微分法及其应用}

\setcounter{section}{5}

\section{多元函数微分学的几何应用}

\subsection{空间曲线的切线与法平面}

{\hei 曲线为参数方程形式} 设曲线 $\Gamma$ 的参数方程为：
\[
\begin{cases}
x = \varphi(t) \\
y = \psi(t) \\
z = \omega(t)
\end{cases}, \quad t \in [\alpha, \beta]
\]
其中 $\varphi, \psi, \omega$ 可导且导数不同时为零。

\begin{itemize}
    \item \textbf{切向量}：$\vec{T} = (\varphi'(t_0), \psi'(t_0), \omega'(t_0))$
    \item \textbf{切线方程}：
    \[
    \frac{x - x_0}{\varphi'(t_0)} = \frac{y - y_0}{\psi'(t_0)} = \frac{z - z_0}{\omega'(t_0)}
    \]
    \item \textbf{法平面方程}：
    \[
    \varphi'(t_0)(x - x_0) + \psi'(t_0)(y - y_0) + \omega'(t_0)(z - z_0) = 0
    \]
\end{itemize}

{\hei 曲线为一般方程形式} 设曲线 $\Gamma$ 由方程组：
\[
\begin{cases}
F(x, y, z) = 0 \\
G(x, y, z) = 0
\end{cases}
\]
确定，$F, G$ 具有连续偏导数，且：
\[
\frac{\partial(F, G)}{\partial(y, z)} \neq 0
\]

\begin{itemize}
    \item \textbf{切向量}：
    \[
    \left(
    \begin{vmatrix}
    F'_y & F'_z \\
    G'_y & G'_z
    \end{vmatrix},
    \begin{vmatrix}
    F'_z & F'_x \\
    G'_z & G'_x
    \end{vmatrix},
    \begin{vmatrix}
    F'_x & F'_y \\
    G'_x & G'_y
    \end{vmatrix}
    \right)
    \]
    
    \item \textbf{切线方程}：
    \[
    \frac{x - x_0}{\left| \begin{matrix} F'_y & F'_z \\ G'_y & G'_z \end{matrix} \right|} = 
    \frac{y - y_0}{\left| \begin{matrix} F'_z & F'_x \\ G'_z & G'_x \end{matrix} \right|} = 
    \frac{z - z_0}{\left| \begin{matrix} F'_x & F'_y \\ G'_x & G'_y \end{matrix} \right|}
    \]
    
    \item \textbf{法平面方程}：
    \[
    \begin{vmatrix}
    F'_y & F'_z \\
    G'_y & G'_z
    \end{vmatrix} (x - x_0) +
    \begin{vmatrix}
    F'_z & F'_x \\
    G'_z & G'_x
    \end{vmatrix} (y - y_0) +
    \begin{vmatrix}
    F'_x & F'_y \\
    G'_x & G'_y
    \end{vmatrix} (z - z_0) = 0
    \]
\end{itemize}

\subsection{曲面的切平面与法线}

{\hei 曲面为显式方程形式}

设曲面方程为 $z = f(x, y)$，$f$ 在点 $(x_0, y_0)$ 处可微。

\begin{itemize}
    \item \textbf{法向量}：$\vec{n} = (f_x'(x_0, y_0), f_y'(x_0, y_0), -1)$
    \item \textbf{切平面方程}：
    \[
    z - z_0 = f_x'(x_0, y_0)(x - x_0) + f_y'(x_0, y_0)(y - y_0)
    \]
    \item \textbf{法线方程}：
    \[
    \frac{x - x_0}{f_x'(x_0, y_0)} = \frac{y - y_0}{f_y'(x_0, y_0)} = \frac{z - z_0}{-1}
    \]
\end{itemize}

{\hei 曲面为隐式方程形式}

设曲面方程为 $F(x, y, z) = 0$，且 $F$ 具有连续偏导数，$(F_x, F_y, F_z) \neq \vec{0}$。

\begin{itemize}
    \item \textbf{法向量}：$\vec{n} = (F_x'(x_0, y_0, z_0), F_y'(x_0, y_0, z_0), F_z'(x_0, y_0, z_0))$
    \item \textbf{切平面方程}：
    \[
    F_x'(x_0, y_0, z_0)(x - x_0) + F_y'(x_0, y_0, z_0)(y - y_0) + F_z'(x_0, y_0, z_0)(z - z_0) = 0
    \]
    \item \textbf{法线方程}：
    \[
    \frac{x - x_0}{F_x'(x_0, y_0, z_0)} = \frac{y - y_0}{F_y'(x_0, y_0, z_0)} = \frac{z - z_0}{F_z'(x_0, y_0, z_0)}
    \]
\end{itemize}

\subsection{全微分的几何意义}

对于二元函数 $z = f(x, y)$：

\begin{itemize}
    \item $\Delta z$：函数在点 $(x_0, y_0)$ 处的实际增量。
    \item $dz$：函数在点 $(x_0, y_0)$ 处的全微分，等于切平面上对应点的竖坐标增量。
    \item \textbf{几何意义}：$dz$ 是切平面上的增量，$\Delta z - dz$ 是比 $\sqrt{\Delta x^2 + \Delta y^2}$ 高阶的无穷小。
\end{itemize}

\[
dz = f_x'(x_0, y_0)\Delta x + f_y'(x_0, y_0)\Delta y
\]

\subsection{法向量的方向余弦（显式曲面）}

对于曲面 $z = f(x, y)$，取向上方向的法向量：

\[
\cos \alpha = \frac{-f_x}{\sqrt{1 + f_x^2 + f_y^2}}, \quad
\cos \beta = \frac{-f_y}{\sqrt{1 + f_x^2 + f_y^2}}, \quad
\cos \gamma = \frac{1}{\sqrt{1 + f_x^2 + f_y^2}}
\]

其中 $\gamma$ 为法向量与 $z$ 轴正方向的夹角。








\end{document}
 